
\begin{table}[H]
\centering
\caption{Phân công nội dung của các thành viên}
\label{tab:app:contributions}
\small
\setlength{\tabcolsep}{8pt}
\renewcommand{\arraystretch}{1.2}
\begin{tabular}{lll}
\toprule
Thành viên & MSSV & Mục phụ trách \\
\midrule
Nguyễn Thị Mai Quyên & 2510213 & \begin{tabular}[t]{@{}l@{}}Chương~1: Đặt vấn đề \\ Mục 2.1--2.4: Giới thiệu bộ dữ liệu NSL-KDD \\ Mục 3.1--3.2: Thiết kế phân tích; thống kê mô tả và trực quan hóa \\ Mục 3.3: Phân tích tương quan Pearson \\ Mục 4.3: Kết quả phân tích tương quan Pearson\end{tabular} \\
\cmidrule(lr){1-3}
Nguyễn Thị Mỹ Uyên & 2510264 & \begin{tabular}[t]{@{}l@{}}Mục 2.5--2.7: Tiền xử lý và quy ước dữ liệu \\ Mục 3.3: Phân tích tương quan Pearson \\ Mục 4.1--4.2: Thống kê mô tả, phân bố nhãn \\ Mục 4.3: Kết quả phân tích tương quan Pearson\end{tabular} \\
\cmidrule(lr){1-3}
Nguyễn Kim Ngọc Trâm & 2510245 & \begin{tabular}[t]{@{}l@{}}Mục 3.4: Khoảng tin cậy Wilson cho tỷ lệ \\ Mục 3.5: ANOVA một yếu tố và phương án Welch \\ Mục 3.6: Kiểm định hai mẫu độc lập \\ Mục 4.4: Kết quả khoảng tin cậy Wilson 95\% \\ Mục 4.5: Kết quả ANOVA, giả định và hậu kiểm \\ Mục 4.6: Kết quả kiểm định hai mẫu độc lập\end{tabular} \\
\cmidrule(lr){1-3}
Trần Minh Khoa & 2510127 & \begin{tabular}[t]{@{}l@{}}Mục 3.7: Kiểm định Chi-square và Cramér's $V$ \\ Mục 4.7: Kết quả kiểm định Chi-square \\ Mục 4.9: Tổng hợp và đối chiếu các phát hiện \\ Mục 4.10: Hạn chế và khả năng khái quát \\ Chương~5: Kết luận và kiến nghị\end{tabular} \\
\cmidrule(lr){1-3}
Nguyễn Ngọc Đô & 2510043 & \begin{tabular}[t]{@{}l@{}}Mục 3.7: Kiểm định Chi-square và Cramér's $V$ \\ Mục 3.8: Hồi quy Logistic và đánh giá mô hình \\ Mục 4.7: Kết quả kiểm định Chi-square \\ Mục 4.8: Kết quả hồi quy Logistic\end{tabular} \\
\bottomrule
\end{tabular}
\end{table}